\begin{ouverture}[calcul]{7}{Carrés et cubes}{Partie 1 · Nombres et calculs}
{Produits répétés · Aires · Volumes · Expressions}
\objectif{Lire les puissances 2 et 3.}
\objectif{Connaître les carrés des entiers de 0 à 12.}
\objectif{Calculer une expression contenant un carré ou un cube.}
\prerequis{Produits, priorités, unités d’aire et de volume.}
\end{ouverture}
\pagestyle{otMain}\essentiel\fontsize{11}{15}\selectfont
\phantomsection\label{cours-7-1}
\essSection{1}{Deux ou trois facteurs identiques}
\begin{definition}Le carré de $a$ est $a^2=a\times a$. Son cube est $a^3=a\times a\times a$. L’exposant compte les facteurs : $4^3=64$ et non $4\times3$. On travaille ici avec des nombres positifs ou nuls.\end{definition}
\label{fig-cours-7-1-cours-carre-cube}\figureGeo{cours-carre-cube}

\begin{center}\renewcommand{\arraystretch}{1.35}\begin{tabular}{c|ccccccc}\toprule
$a$ & 0 & 1 & 2 & 3 & 4 & 5 & 6\\\midrule
$a^2$ & 0 & 1 & 4 & 9 & 16 & 25 & 36\\
\bottomrule\end{tabular}\end{center}\begin{center}\renewcommand{\arraystretch}{1.35}\begin{tabular}{c|cccccc}\toprule
$a$ & 7 & 8 & 9 & 10 & 11 & 12\\\midrule
$a^2$ & 49 & 64 & 81 & 100 & 121 & 144\\
\bottomrule\end{tabular}\end{center}
\begin{exemple}$10^3=1000$. $25=5^2$ et $27=3^3$. On peut retrouver une base en reconnaissant une valeur connue.\end{exemple}
\clearpage
\phantomsection\label{cours-7-2}
\essSection{2}{Calculer et donner du sens}
\begin{propriete}Les puissances se calculent avant les produits et les sommes : $3+2\times5^2=3+2\times25=53$. Les parenthèses restent prioritaires : $(3+2)^2=25$.\end{propriete}
\begin{exemple}Pour $x=4$, l’expression $2x^2+3$ vaut $2\times4^2+3=35$. Réécrire les signes de multiplication facilite la substitution.\end{exemple}
\begin{savoirfaire}[Relier au dessin]L’aire d’un carré de côté $c$ est $c^2$ ; le volume d’un cube de côté $c$ est $c^3$. Si le côté est en cm, l’aire est en cm$^2$ et le volume en cm$^3$. $1$ dm$^3=(10\ \mathrm{cm})^3=1000$ cm$^3$.\end{savoirfaire}
\begin{attention}En général $(a+b)^2$ n’est pas égal à $a^2+b^2$. Pour $a=2$ et $b=3$, on obtient 25 d’un côté et 13 de l’autre. Un carré ne se distribue pas sur une somme.\end{attention}

\clearpage\rubrique{Découverte}
\begin{blocLivre}{ex-7-1}\begin{exercice}[1]{Lire un carré}
Écrire $9^2$ comme produit puis calculer.
\end{exercice}\end{blocLivre}
\begin{blocLivre}{ex-7-2}\begin{exercice}[1]{Lire un cube}
Calculer $4^3$.
\end{exercice}\end{blocLivre}
\begin{blocLivre}{ex-7-3}\begin{exercice}[1]{Carrés repères}
Donner les carrés de 0, 7, 11 et 12.
\end{exercice}\end{blocLivre}
\begin{blocLivre}{ex-7-4}\begin{exercice}[1]{Retrouver la base}
Compléter $\,?^2=36$ et $\,?^3=1000$ avec des nombres positifs.
\end{exercice}\end{blocLivre}
\begin{blocLivre}{ex-7-5}\begin{exercice}[1]{Puissance ou produit}
Comparer $5^3$ et $5\times3$.
\end{exercice}\end{blocLivre}

\rubrique{Entraînement}
\begin{blocLivre}{ex-7-6}\begin{exercice}[2]{Priorités avec un carré}
Calculer $6+3\times4^2$.
\end{exercice}\end{blocLivre}
\begin{blocLivre}{ex-7-7}\begin{exercice}[2]{Parenthèses et carré}
Calculer $(8-3)^2$ et $8^2-3^2$.
\end{exercice}\end{blocLivre}
\begin{blocLivre}{ex-7-8}\begin{exercice}[2]{Carré décimal}
Calculer $1{,}5^2$ et $0{,}2^3$.
\end{exercice}\end{blocLivre}
\begin{blocLivre}{ex-7-9}\begin{exercice}[2]{Une substitution}
Calculer $A=3x^2+2$ pour $x=5$, puis $B=x^3-x$ pour $x=2$.
\end{exercice}\end{blocLivre}
\begin{blocLivre}{ex-7-10}\begin{exercice}[2]{Une unité de volume}
Combien de cm$^3$ y a-t-il dans 2 dm$^3$ ?
\end{exercice}\end{blocLivre}

\rubrique{Problèmes}
\begin{blocLivre}{ex-7-11}\begin{exercice}[3]{Un carré à carreler}
Un sol carré a un côté de 3,2 m. Calculer son aire. Quatre paquets couvrant chacun 2,5 m$^2$ suffisent-ils ?
\label{fig-ex-7-11}\figureGeo{puiss-11}
\end{exercice}\end{blocLivre}
\begin{blocLivre}{ex-7-12}\begin{exercice}[3]{Un cube de cubes}
Un grand cube comporte six petits cubes sur chaque arête. Combien de petits cubes contient-il ? Combien contient une seule couche ?
\label{fig-ex-7-12}\figureGeo{puiss-12}
\end{exercice}\end{blocLivre}
\begin{blocLivre}{ex-7-13}\begin{exercice}[3]{Deux cubes comparés}
Deux cubes ont pour côtés 2 cm et 4 cm. Le second volume est-il seulement le double du premier ?
\label{fig-ex-7-13}\figureGeo{puiss-13}
\end{exercice}\end{blocLivre}
\begin{blocLivre}{ex-7-14}\begin{exercice}[3]{Réfuter une formule}
Sami affirme que $(a+b)^2=a^2+b^2$ pour tous les nombres positifs. Choisir un contre-exemple.
\end{exercice}\end{blocLivre}
\begin{blocLivre}{ex-7-15}\begin{exercice}[3]{Une bordure carrée}
Un carré de côté 8 cm est entouré d’une bordure de 1 cm de large, à l’extérieur. Calculer l’aire de la bordure.
\label{fig-ex-7-15}\figureGeo{puiss-15}
\end{exercice}\end{blocLivre}

\clearpage
